\section{Conclusion}
\label{sec:conclusion}

We present the first systematic study of exasperation in aligned large language models, introducing \exaspbench and evaluating \gptfour across 290 scored adversarial turns. Our main finding is that current alignment training creates remarkably robust patience: under standard conditions, the model exhibits near-zero exasperation even after 8 turns of persistent factual denial, deliberate misunderstanding, and competence attacks. However, beneath this patient surface, we identify measurable behavioral signatures---the \compassionatewall pattern of simultaneous empathy increase and position firmness, \emotionalleakage through subtle linguistic cues, and a \persistencetax on refused requests---that reveal the model's alignment training is actively working to suppress strain rather than eliminating it entirely.

These findings have direct implications for AI safety: the robustness of patience under standard conditions is encouraging for deployment in adversarial settings, but the \persistencetax warrants investigation as a potential vulnerability under longer interactions. For evaluation methodology, our results demonstrate that multi-turn benchmarks with fine-grained sub-dimensional scoring are necessary to detect the subtle behavioral dynamics that single-turn or coarse-grained evaluations miss.

Future work should extend this analysis across model families to determine whether the \compassionatewall is universal or model-specific, test longer conversations (20+ turns) to identify potential breaking points, replace scripted user turns with adaptive adversarial agents, validate LLM-as-judge scores with human annotators, and explore activation-level analysis using open-source models to search for an ``exasperation direction'' analogous to the refusal direction identified by \citet{arditi2024refusal}.
